\section{Planar equivalent layers}
\label{eq:planar_eqlayer}

Consider a continuous equivalent layer $p_{\varepsilon}(\mathbf{r}'')$ (Definition \partoneref{def:continuous-eqlayer})  
defined on a horizontal surface $\mathcal{S}''$, which may either lie above the sources or intersect them  
(Figure \partoneref{fig:Figure1}). In practical applications, however, a continuous layer is not directly used.  
Instead, one typically estimates a \textit{discrete equivalent layer} that is sufficiently smooth and fits the observed  
discrete data within a prescribed error tolerance (Theorem \partoneref{theorem:planar-eqlayer-theorem}, Remark 2). 
We regard any discrete layer commonly obtained in practice  
as an approximation of a corresponding continuous equivalent layer $p_{\varepsilon}(\mathbf{r}'')$  
(Definition \partoneref{def:continuous-eqlayer}). A \textit{discrete equivalent layer} on the horizontal surface  
$\mathcal{S}''$ (Figure \partoneref{fig:Figure1}) can be formally defined as follows:  

\begin{definition}[\textbf{Discrete equivalent layer}]
	Let us assume to have an $N \times 1$ observed data vector $\mathbf{d}^{o}$ whose elements 
	are the potential-field data $d^{o}_{i}$ at observation points $\mathbf{r}_{i}$ outside the sources, 
	$i \in \{1, \dots, N\}$, above a horizontal surface $\mathcal{S}''$ with constant vertical 
	coordinate $z''_{0}$. Besides, consider an $N \times M$ 
	matrix $\mathbf{A}$ with element $a_{ij}$ defined by a harmonic function that depends only on $\mathbf{r}_{i}$ 
	and $\mathbf{r}''_{j} \in \mathcal{S}''$, $j \in \{1, \dots, M\}$, and
	vanishes as the distance $\| \mathbf{r}_{i} - \mathbf{r}''_{j} \|_{2}$ increases. 
	Then, for a chosen accuracy $\epsilon > 0$, we may call the $M \times 1$ vector $\mathbf{p}_{\epsilon}$ with $j$-th element $p_{j}$ 
	a \underline{discrete equivalent layer} if its predicted potential field 
	\begin{equation}
		\mathbf{d}_{\epsilon} = \mathbf{A} \, \mathbf{p}_{\epsilon}
		\label{eq:predicted-data-discrete-layer}
	\end{equation}
	satisfies 
	\begin{equation}
		\| \mathbf{d}^{o} - \mathbf{d}_{\epsilon} \|_{2}^{2} \le \epsilon \: .
		\label{eq:epsilon-bound-data-discrete}
	\end{equation}
	\label{def:discrete-eqlayer}
\end{definition}

In Definition \ref{def:discrete-eqlayer}, $\mathbf{d}_{\epsilon}$ is the predicted data vector whose 
$i$-th element represents the potential field due to $\mathbf{p}_{\epsilon}$ at the observation point $\mathbf{r}_{i}$. 
Besides, the elements of $\mathbf{A}$ are usually defined by 
$a_{ij} \equiv D_{a} \Psi(\mathbf{r}_{i}, \mathbf{r}''_{j})$, which represents a 
partial derivative of $\Psi$ (equation \partoneref{eq:Psi}), with $D_{a}$ given by the space-domain differential operator 
$D$ (equation \partoneref{eq:diff-operators-generic}). 
It is important to bear in mind that the horizontal surface $\mathcal{S}''$ in Definition \ref{def:discrete-eqlayer} 
may cross the sources, as illustrated in Figure \partoneref{fig:Figure1}.
The main assumption in Definition \ref{def:discrete-eqlayer} is that $d^{o}_{i}$ can be approximated 
within a given accuracy $\epsilon$ by the potential field produced by a discrete and finite set of $M$ 
virtual sources such as monopoles (or point masses) or dipoles with 
constant and unknown physical property $p_{j}$ (e.g., mass, magnetic-moment intensity). 
These virtual sources are usually called \textit{equivalent sources}.
We emphasize that Definition \ref{def:discrete-eqlayer} presumes a planar equivalent layer formed 
by punctual equivalent sources that satisfies the data fit criterion given by equation \ref{eq:epsilon-bound-data-discrete}, 
but these premises are certainly open to adjustments if necessary.

After estimating a particular $\mathbf{p}_{\epsilon}$ (Definition \ref{def:discrete-eqlayer}), 
a potential field transformation such as upward continuation, 
interpolation or reduction to the pole, for example, can be computed as follows:
\begin{equation}
	\mathbf{f}_{\epsilon} = \mathbf{T} \, \mathbf{p}_{\epsilon} \: ,
	\label{eq:transformed-field-discrete-layer} 
\end{equation}
where $\mathbf{f}_{\epsilon}$ is an $L \times 1$ vector whose $k$-th element approximates 
the transformed potential field $f(\mathbf{r}^{\ddagger}_{k})$ (equation \partoneref{eq:transformed-field-space}) 
at a point $\mathbf{r}^{\ddagger}_{k} \in \mathcal{S}^{\ddagger}$ above $\mathcal{S}''$ (Figure \partoneref{fig:Figure1})
and $\mathbf{T}$ is an $L \times M$ matrix whose element $t_{kj} \equiv D_{t} \Psi(\mathbf{r}^{\ddagger}_{k}, \mathbf{r}''_{j})$,
$k \in \{ 1, \dots, L \}$, $j \in \{ 1, \dots, M \}$, 
is commonly defined by a partial derivative of $\Psi$ (equation \partoneref{eq:Psi}), with $D_{t}$ given by the 
space-domain differential operator $D$ (equation \partoneref{eq:diff-operators-generic}).
The element $t_{kj}$ represents the transformed potential field due to the $j$-th equivalent source. 
We have defined the transformed potential field at points 
$\mathbf{r}^{\ddagger}_{k} \in \mathcal{S}^{\ddagger}$ above $\mathcal{S}''$ (Figure \partoneref{fig:Figure1}), however they 
can be defined at any place above $\mathcal{S}''$, including the observation points. 

From the physical point of view, the product of the $j$-th element $p_{j}$ of $\mathbf{p}_{\epsilon}$ 
and the harmonic function $a_{ij}$ representing the element $ij$ of matrix $\mathbf{A}$ 
(equation \ref{eq:predicted-data-discrete-layer}) defines the potential field produced at the $i$-th observation
point by the $j$-th equivalent source with scalar physical property $p_{j}$.
For example, if $\mathbf{d}^{o}$ contains gravity disturbance data
(Table \partoneref{tbl:Table2}), then $p_{j} \, a_{ij}$ commonly defines the vertical component of the gravitational attraction
produced at $\mathbf{r}_{i}$ by a monopole located at
$\mathbf{r}''_{j} \in \mathcal{S}''$ (Figure \partoneref{fig:Figure1}), with mass $p_{j}$. 
This is the pioneer approach presented by \cite{dampney1969} for 
processing gravity data.
For the case in which $\mathbf{d}^{o}$ contains total-field anomaly data 
\cite[e.g.,][ p. 178]{blakely1996} instead, the common approach consists in
following \cite{emilia1973} and defining the term $p_{j} \, a_{ij}$ such
that it represents the approximated total-field anomaly (Table \partoneref{tbl:Table2}) 
produced at the $i$-th observation point by a dipole located at 
$\mathbf{r}''_{j} \in \mathcal{S}''$ (Figure \partoneref{fig:Figure1}), with magnetic moment intensity $p_{j}$.

Note that, in the examples presented above, the physical relationship 
between the observed potential-field data and their true sources is the
same as that between the data and the equivalent sources.
However, many authors 
\cite[e.g.,][]{cordell1992, guspi_novara2009, li_etal2014, soler_uieda2021} 
have proposed equivalent-layer methods by assuming without any proof that the harmonic functions $a_{ij}$ are 
independent of the nature of the observed potential-field data and can be 
conveniently chosen according to the characteristics of the problem.
We mathematically confirm this independence some sections later and also show how 
the particular form of predefined harmonic functions $a_{ij}$ affect the subsequent 
potential-field transformations.